\part[3]\hfill

  Consider the following declarations, evaluated in order:

  \begin{codeblock}
    fun makeAdder () =
      let
        val n = 42
      in
        fn y => n + y
      end

    val add = makeAdder ()
    val n = 1000
    val res = add 1
  \end{codeblock}

  What is the value of \code{res}? Justify your answer in one sentence.

  \begin{solutionorbox}[9em]
    \code{res} $\eeq$ \code{43}.

    The function \code{fn y => n + y} captures the binding of \code{n} in scope
    \textbf{where it was defined} --- the \code{val n = 42} inside
    \code{makeAdder} --- not whatever \code{n} happens to be bound to at the
    point where \code{add} is later called. The later \code{val n = 1000} is a
    different, unrelated binding and has no effect on \code{add}.
  \end{solutionorbox}


% -----------------------------------------------------------------------------
% 3. Equivalence is not runtime cost.
%    Misunderstanding: students conflate "provably equivalent" with "the same
%    program" / "same efficiency", and think ~= says anything about cost.
% -----------------------------------------------------------------------------
